\documentclass[12pt,a4paper]{article}
\usepackage[utf8]{inputenc}
\usepackage[margin=1in]{geometry}
\usepackage{hyperref}
\usepackage{graphicx}
\usepackage{enumitem}
\usepackage{fancyhdr}
\usepackage{titlesec}
\usepackage{xcolor}
\usepackage{tcolorbox}
\usepackage{tabularx}
\usepackage{longtable}

% Define colors
\definecolor{primarycolor}{RGB}{0,102,204}
\definecolor{secondarycolor}{RGB}{51,51,51}

% Header and footer
\pagestyle{fancy}
\fancyhf{}
\fancyhead[L]{CSEDU Digital Knowledge Platform}
\fancyhead[R]{User Manual}
\fancyfoot[C]{\thepage}

% Title formatting
\titleformat{\section}{\Large\bfseries\color{primarycolor}}{\thesection}{1em}{}
\titleformat{\subsection}{\large\bfseries\color{secondarycolor}}{\thesubsection}{1em}{}

% Hyperlink colors
\hypersetup{
    colorlinks=true,
    linkcolor=primarycolor,
    urlcolor=primarycolor,
    citecolor=primarycolor
}

\title{\textbf{\Huge CSEDU Digital Knowledge Platform}\\
\vspace{0.5cm}
\Large Comprehensive User Manual}
\author{Department of Computer Science and Engineering\\
University of Dhaka}
\date{\today}

\begin{document}

\maketitle
\thispagestyle{empty}
\newpage

\tableofcontents
\newpage

\section{Introduction}

\subsection{Overview}
The CSEDU Digital Knowledge Platform is a comprehensive web-based system designed to facilitate knowledge sharing, research collaboration, and library management within the Computer Science and Engineering Department at the University of Dhaka. The platform integrates multiple services including a digital library, research repository, student project showcase, multimedia archives, and AI-powered assistance.

\subsection{Key Features}
\begin{itemize}[leftmargin=*]
    \item \textbf{Digital Library Management} - Browse, borrow, and return books with automated fine calculations
    \item \textbf{Research Repository} - Submit, review, and publish academic research papers
    \item \textbf{Project Showcase} - Display and archive student projects with team collaborations
    \item \textbf{Digital Archives} - Store and access historical documents, multimedia content
    \item \textbf{AI Assistant} - Intelligent search and document summarization using RAG technology
    \item \textbf{Role-Based Access Control} - Five-tier permission system ensuring appropriate access
    \item \textbf{Notification System} - Email and in-app notifications for important events
\end{itemize}

\subsection{User Roles}
The platform supports five distinct user roles, each with specific permissions:

\begin{enumerate}[leftmargin=*]
    \item \textbf{Public (Unauthenticated)} - View published content only
    \item \textbf{Student} - Default role with basic borrowing and upload permissions
    \item \textbf{Researcher} - Enhanced access to submit and review research papers
    \item \textbf{Librarian} - Manage library catalog, approve content, process loans
    \item \textbf{Administrator} - Full platform management and user role administration
\end{enumerate}

\subsection{Accessing the Platform}
\begin{tcolorbox}[colback=blue!5!white,colframe=primarycolor,title=Platform Access]
\textbf{URL:} \url{http://20.195.127.226:8080}\\
\textbf{Support:} For assistance, contact the CSE Department IT support team.
\end{tcolorbox}

\subsection{System Requirements}
\begin{itemize}[leftmargin=*]
    \item Modern web browser (Chrome, Firefox, Safari, Edge - latest versions)
    \item Internet connection
    \item Recommended screen resolution: 1366x768 or higher
    \item JavaScript enabled
\end{itemize}

\newpage

\section{Getting Started}

\subsection{Creating an Account}

\subsubsection{Registration Process}
\begin{enumerate}[leftmargin=*]
    \item Navigate to the platform homepage
    \item Click on the \textbf{"Sign Up"} or \textbf{"Register"} button
    \item Fill in the registration form:
        \begin{itemize}
            \item \textbf{Full Name:} Your complete name as per university records
            \item \textbf{Email Address:} A valid email address (preferably institutional email)
            \item \textbf{Password:} Minimum 8 characters, include uppercase, lowercase, numbers, and special characters
            \item \textbf{Confirm Password:} Re-enter your password
        \end{itemize}
    \item Click \textbf{"Create Account"}
    \item Check your email for a confirmation message
    \item Your account is created with \textbf{Student} role by default
\end{enumerate}

\begin{tcolorbox}[colback=yellow!10!white,colframe=orange,title=Important Note]
All new registrations are assigned the \textbf{Student} role automatically. To request elevated permissions (Researcher or Librarian), use the Role Request feature after logging in.
\end{tcolorbox}

\subsection{Logging In}
\begin{enumerate}[leftmargin=*]
    \item Click the \textbf{"Login"} button on the homepage
    \item Enter your registered email address
    \item Enter your password
    \item Click \textbf{"Sign In"}
    \item You will be redirected to your personalized dashboard
\end{enumerate}

\subsection{Password Recovery}
If you forget your password:
\begin{enumerate}[leftmargin=*]
    \item Click \textbf{"Forgot Password?"} on the login page
    \item Enter your registered email address
    \item Check your email for a password reset link
    \item Click the link (valid for 1 hour)
    \item Create a new password (minimum 8 characters)
    \item Confirm your new password
    \item Click \textbf{"Reset Password"}
    \item Log in with your new credentials
\end{enumerate}

\subsection{Profile Management}

\subsubsection{Viewing Your Profile}
\begin{enumerate}[leftmargin=*]
    \item Click on your name/avatar in the top-right corner
    \item Select \textbf{"Profile"} from the dropdown menu
    \item View your profile information:
        \begin{itemize}
            \item Name, Email, Role
            \item Registration date
            \item Last login time
            \item Account statistics (uploads, loans, etc.)
        \end{itemize}
\end{enumerate}

\subsubsection{Editing Your Profile}
\begin{enumerate}[leftmargin=*]
    \item Navigate to your profile page
    \item Click the \textbf{"Edit Profile"} button
    \item Update your information (Name, Email)
    \item Click \textbf{"Save Changes"}
    \item Confirmation message will appear
\end{enumerate}

\subsubsection{Changing Your Password}
\begin{enumerate}[leftmargin=*]
    \item Go to \textbf{Settings} or \textbf{Profile}
    \item Click \textbf{"Change Password"}
    \item Enter your current password
    \item Enter your new password (minimum 8 characters)
    \item Confirm your new password
    \item Click \textbf{"Update Password"}
    \item You will be logged out and prompted to log in again
\end{enumerate}

\newpage

\section{Library Catalog System}

\subsection{Browsing the Catalog}
The library catalog is accessible to all users, including non-authenticated visitors.

\subsubsection{Accessing the Catalog}
\begin{enumerate}[leftmargin=*]
    \item Click \textbf{"Library"} or \textbf{"Catalog"} in the main navigation menu
    \item The catalog page displays all available books in a grid or list view
    \item Each book card shows:
        \begin{itemize}
            \item Book title and author
            \item Cover image (if available)
            \item Availability status (Available/Borrowed)
            \item Format (Book, E-book, etc.)
            \item Year of publication
        \end{itemize}
\end{enumerate}

\subsubsection{Searching for Books}
\begin{enumerate}[leftmargin=*]
    \item Use the search bar at the top of the catalog page
    \item Enter keywords from:
        \begin{itemize}
            \item Book title
            \item Author name
            \item ISBN number
            \item Topic/subject area
        \end{itemize}
    \item Press Enter or click the search icon
    \item Results are displayed with relevance ranking
\end{enumerate}

\subsubsection{Filtering Books}
You can narrow down catalog results using multiple filters:
\begin{itemize}[leftmargin=*]
    \item \textbf{Topic:} Filter by subject area (e.g., Algorithms, Machine Learning, Databases, Programming)
    \item \textbf{Availability:} Show only available or borrowed books
    \item \textbf{Format:} Book, E-book, Reference
    \item \textbf{Year:} Publication year range
\end{itemize}

\subsection{Borrowing Books (Student, Researcher, Librarian, Administrator)}

\subsubsection{How to Borrow}
\begin{enumerate}[leftmargin=*]
    \item Log in to your account
    \item Navigate to the library catalog
    \item Find the book you want to borrow (must show "Available")
    \item Click on the book to view details
    \item Click the \textbf{"Borrow"} button
    \item Confirm the borrowing action
    \item A success message appears with the due date
    \item The book is now listed in \textbf{"My Loans"}
\end{enumerate}

\subsubsection{Loan Period and Limits}
\begin{itemize}[leftmargin=*]
    \item \textbf{Loan Duration:} 14 days from the checkout date
    \item \textbf{Due Date:} Displayed in your loans list
    \item \textbf{Late Returns:} Subject to fines (50 BDT per day, max 500 BDT per loan)
    \item \textbf{Fine Block:} If outstanding fines exceed 50 BDT, borrowing is disabled until payment
\end{itemize}

\begin{tcolorbox}[colback=red!5!white,colframe=red,title=Borrowing Restrictions]
\begin{itemize}[leftmargin=*]
    \item You cannot borrow the same book twice simultaneously
    \item Librarians cannot borrow books (monitoring role only)
    \item Borrowing is disabled if you have overdue fines exceeding 50 BDT
\end{itemize}
\end{tcolorbox}

\subsection{Viewing Your Loans}
\begin{enumerate}[leftmargin=*]
    \item Click \textbf{"My Loans"} in the navigation menu
    \item See all your current and past loans
    \item For each loan, view:
        \begin{itemize}
            \item Book title
            \item Checkout date
            \item Due date
            \item Return date (if returned)
            \item Status (Active, Returned, Overdue)
        \end{itemize}
    \item Active loans show a \textbf{"Return"} button
\end{enumerate}

\subsection{Returning Books}
\begin{enumerate}[leftmargin=*]
    \item Go to \textbf{"My Loans"}
    \item Find the book you want to return
    \item Click the \textbf{"Return Book"} button
    \item Confirm the return
    \item If the book is overdue, a fine will be automatically calculated
    \item The book status changes to "Returned"
    \item Book availability is updated in the catalog
\end{enumerate}

\newpage

\subsection{Reservations (Holds) System}

\subsubsection{Placing a Hold}
When all copies of a book are borrowed, you can place a hold to be notified when it becomes available.

\begin{enumerate}[leftmargin=*]
    \item Navigate to the book detail page
    \item If all copies are borrowed, a \textbf{"Place Hold"} button appears
    \item Click \textbf{"Place Hold"}
    \item Your hold request is recorded with a timestamp
    \item You'll receive an email notification when the book becomes available
    \item The hold is valid for 48 hours after notification
\end{enumerate}

\subsubsection{Viewing Your Holds}
\begin{enumerate}[leftmargin=*]
    \item Click \textbf{"My Loans"} or \textbf{"Holds"} in the navigation
    \item View all your active reservations
    \item See your position in the queue
    \item Status indicators:
        \begin{itemize}
            \item \textbf{Active:} Waiting for book availability
            \item \textbf{Fulfilled:} Book is now available for you
            \item \textbf{Cancelled:} You or staff cancelled the hold
        \end{itemize}
\end{enumerate}

\subsubsection{Cancelling a Hold}
\begin{enumerate}[leftmargin=*]
    \item Go to your holds list
    \item Find the hold you want to cancel
    \item Click the \textbf{"Cancel Hold"} button
    \item Confirm cancellation
    \item The hold is removed from your list
\end{enumerate}

\subsection{Fine Management}

\subsubsection{Viewing Your Fines (Student, Researcher)}
\begin{enumerate}[leftmargin=*]
    \item Click \textbf{"Fines"} in the navigation menu
    \item See all your fines with details:
        \begin{itemize}
            \item Book title
            \item Fine amount (BDT)
            \item Status (Pending, Paid, Waived)
            \item Calculation date
        \end{itemize}
    \item Total unpaid fines are displayed at the top
\end{enumerate}

\subsubsection{Fine Calculation}
\begin{itemize}[leftmargin=*]
    \item \textbf{Rate:} 50 BDT per day overdue
    \item \textbf{Maximum:} 500 BDT per loan
    \item \textbf{Calculation:} Automatic on the day after due date
    \item \textbf{Grace Period:} None - fines start immediately after the due date
\end{itemize}

\subsubsection{Paying Fines}
\begin{enumerate}[leftmargin=*]
    \item Navigate to \textbf{"Fines"} page
    \item Find the pending fine you want to pay
    \item Click \textbf{"Pay Fine"}
    \item Select payment method (if configured)
    \item Confirm payment
    \item Fine status changes to "Paid"
    \item If total fines were blocking borrowing, the block is lifted
\end{enumerate}

\begin{tcolorbox}[colback=yellow!10!white,colframe=orange,title=Payment Note]
Contact the library desk or administrator for payment procedures. The system records payment confirmations entered by library staff.
\end{tcolorbox}

\newpage

\section{Digital Archives}

\subsection{Overview}
The Digital Archives section stores historical documents, multimedia content, and departmental resources accessible to all users.

\subsection{Browsing Archives}
\begin{enumerate}[leftmargin=*]
    \item Click \textbf{"Archives"} in the main navigation
    \item Browse through categories or search
    \item Each archive item displays:
        \begin{itemize}
            \item Title
            \item Description/Abstract
            \item Upload date
            \item Format (PDF, Video, Image, Document, etc.)
            \item Access tier (Public, Student, Researcher, Restricted)
        \end{itemize}
    \item Click on an item to view full details
\end{enumerate}

\subsection{Uploading to Archives (Student, Researcher, Librarian, Administrator)}

\subsubsection{Supported File Types}
\begin{itemize}[leftmargin=*]
    \item \textbf{Documents:} PDF, DOCX, DOC, PPTX, PPT, XLSX, XLS
    \item \textbf{Images:} JPG, JPEG, PNG, GIF
    \item \textbf{Videos:} MP4
    \item \textbf{Audio:} MP3
    \item \textbf{Archives:} ZIP
    \item \textbf{External Links:} URLs (YouTube videos, web resources)
    \item \textbf{Maximum File Size:} 50 MB per file
\end{itemize}

\subsubsection{Upload Process}
\begin{enumerate}[leftmargin=*]
    \item Click \textbf{"Upload"} in the navigation
    \item Select \textbf{"Archive"} as the upload type
    \item Fill in the form:
        \begin{itemize}
            \item \textbf{Title:} Descriptive title (required)
            \item \textbf{Abstract:} Brief description of the content
            \item \textbf{Keywords:} Comma-separated tags for searchability
            \item \textbf{Tags:} Additional categorization tags
            \item \textbf{Language:} Content language (default: English)
            \item \textbf{Access Tier:} Choose who can view (Public, Student, Researcher, Restricted)
            \item \textbf{Status:} Draft (save for later) or Published (make available immediately)
        \end{itemize}
    \item Choose file source:
        \begin{itemize}
            \item Upload a file from your computer, OR
            \item Provide an external URL (for YouTube videos, web resources)
        \end{itemize}
    \item Click \textbf{"Upload"}
    \item Wait for processing confirmation
    \item Your upload appears in \textbf{"My Uploads"}
\end{enumerate}

\subsection{Viewing Archive Details}
\begin{enumerate}[leftmargin=*]
    \item Click on any archive item
    \item View complete details:
        \begin{itemize}
            \item Title, abstract, keywords
            \item Uploader name and date
            \item File format and size
            \item Access permissions
        \end{itemize}
    \item Options available:
        \begin{itemize}
            \item \textbf{Download:} Save file to your device
            \item \textbf{View Online:} Preview (for PDFs, images, videos)
            \item \textbf{Share:} Copy link to share with others
        \end{itemize}
\end{enumerate}

\subsection{Downloading Archive Files}
\begin{enumerate}[leftmargin=*]
    \item Navigate to the archive item detail page
    \item Click the \textbf{"Download"} button
    \item The file downloads to your default download folder
    \item For external links, you're redirected to the external resource
\end{enumerate}

\subsection{Managing Your Uploads}
\begin{enumerate}[leftmargin=*]
    \item Click \textbf{"My Uploads"} in the navigation
    \item View all your uploaded content (Archives, Research, Projects)
    \item For each item:
        \begin{itemize}
            \item View status (Draft, Review, Published, Archived)
            \item Edit metadata (title, abstract, keywords)
            \item Delete drafts
            \item Submit for review/publication
        \end{itemize}
    \item Click \textbf{"Edit"} to update metadata
    \item Click \textbf{"Delete"} to remove draft items (published items require admin approval)
\end{enumerate}

\newpage

\section{Research Paper Management}

\subsection{Overview (Researcher, Librarian, Administrator Only)}
The Research section allows researchers to submit, review, and publish academic papers, facilitating knowledge sharing and collaboration.

\subsection{Browsing Research Papers (All Users)}
\begin{enumerate}[leftmargin=*]
    \item Click \textbf{"Research"} in the navigation menu
    \item Browse published research papers
    \item Use search and filters:
        \begin{itemize}
            \item Search by title, author, keywords
            \item Filter by publication year
            \item Filter by research area/topic
        \end{itemize}
    \item Each paper displays:
        \begin{itemize}
            \item Title and authors
            \item Publication date
            \item Abstract preview
            \item DOI (if available)
            \item Download option
        \end{itemize}
\end{enumerate}

\subsection{Submitting Research Papers (Researcher, Administrator)}

\subsubsection{Preparing Your Submission}
Before submitting, ensure you have:
\begin{itemize}[leftmargin=*]
    \item Final version of your paper (PDF format, max 50 MB)
    \item Complete list of authors with affiliations
    \item Abstract (required)
    \item Keywords (5-10 recommended)
    \item DOI, journal/conference details (if published)
\end{itemize}

\subsubsection{Submission Process}
\begin{enumerate}[leftmargin=*]
    \item Click \textbf{"Upload"} then select \textbf{"Research Paper"}
    \item Fill in the submission form:
        \begin{itemize}
            \item \textbf{Title:} Full paper title (required)
            \item \textbf{Authors:} Enter all authors' names (comma-separated)
            \item \textbf{Co-authors:} Additional contributors
            \item \textbf{Abstract:} Concise summary of your research (required)
            \item \textbf{Keywords:} Key terms for discoverability
            \item \textbf{DOI:} Digital Object Identifier (if available)
            \item \textbf{Journal/Conference:} Publication venue
            \item \textbf{Publication Date:} When published (if applicable)
            \item \textbf{Access Tier:} Choose visibility (Public, Student, Researcher)
        \end{itemize}
    \item Upload your PDF file
    \item Choose initial status:
        \begin{itemize}
            \item \textbf{Draft:} Save for later editing
            \item \textbf{Submit for Review:} Send to reviewers immediately
        \end{itemize}
    \item Click \textbf{"Submit Research Paper"}
    \item Confirmation message appears
    \item Your submission is listed in \textbf{"My Uploads"}
\end{enumerate}

\subsection{Paper Review Workflow}

\subsubsection{Submission Statuses}
\begin{enumerate}[leftmargin=*]
    \item \textbf{Draft:} Author is still editing
    \item \textbf{Under Review:} Assigned to reviewer for evaluation
    \item \textbf{Published:} Approved and publicly available
    \item \textbf{Archived:} Withdrawn or removed from active display
\end{enumerate}

\subsubsection{Submitting for Review}
\begin{enumerate}[leftmargin=*]
    \item Go to \textbf{"My Uploads"}
    \item Find your draft paper
    \item Click \textbf{"Submit for Review"}
    \item Confirm submission
    \item Status changes to "Under Review"
    \item Librarian/Administrator assigns a reviewer
    \item You'll be notified when review is complete
\end{enumerate}

\subsubsection{Receiving Feedback}
\begin{enumerate}[leftmargin=*]
    \item Check your notifications or email
    \item Navigate to your paper in \textbf{"My Uploads"}
    \item Click to view review comments
    \item Review notes include:
        \begin{itemize}
            \item Reviewer feedback
            \item Suggestions for improvement
            \item Publication decision
        \end{itemize}
    \item If revisions requested:
        \begin{itemize}
            \item Make necessary changes
            \item Re-upload revised version
            \item Resubmit for review
        \end{itemize}
\end{enumerate}

\newpage

\subsection{Viewing Published Papers}
\begin{enumerate}[leftmargin=*]
    \item Click on any published paper
    \item View full details:
        \begin{itemize}
            \item Complete abstract
            \item Author information
            \item Publication details (journal, DOI, date)
            \item Full-text PDF
        \end{itemize}
    \item Download the PDF
    \item Cite the paper (citation format provided)
    \item Use AI to generate summary
\end{enumerate}

\subsection{Editing Your Papers}
\begin{enumerate}[leftmargin=*]
    \item Go to \textbf{"My Uploads"}
    \item Find the paper you want to edit
    \item Click \textbf{"Edit"}
    \item Update metadata (title, abstract, keywords, etc.)
    \item Replace PDF file if needed
    \item Save changes
    \item If already published, changes may require admin approval
\end{enumerate}

\section{Student Projects}

\subsection{Overview (Student, Researcher, Administrator)}
The Projects section showcases student team projects, providing a portfolio of academic and personal work.

\subsection{Browsing Projects (All Users)}
\begin{enumerate}[leftmargin=*]
    \item Click \textbf{"Projects"} in the navigation
    \item Browse the project gallery
    \item Filter by:
        \begin{itemize}
            \item Academic year
            \item Course code
            \item Project type
        \end{itemize}
    \item Each project card shows:
        \begin{itemize}
            \item Project title
            \item Team members
            \item Supervisor
            \item Academic year
            \item Brief description
        \end{itemize}
\end{enumerate}

\subsection{Submitting a Project (Student, Researcher, Administrator)}

\subsubsection{Before Submission}
Prepare the following:
\begin{itemize}[leftmargin=*]
    \item Project title
    \item Team member names (all contributors)
    \item Supervisor name (optional)
    \item Course code and academic year
    \item Project abstract/description
    \item Keywords
    \item Project deliverables:
        \begin{itemize}
            \item Documentation (PDF, DOCX)
            \item Source code (ZIP archive)
            \item Demo video (MP4 or YouTube link)
            \item Web URL (if applicable)
            \item GitHub repository link
            \item APK file (for Android apps)
        \end{itemize}
\end{itemize}

\subsubsection{Submission Steps}
\begin{enumerate}[leftmargin=*]
    \item Click \textbf{"Upload"} then select \textbf{"Project"}
    \item Fill in project details:
        \begin{itemize}
            \item \textbf{Title:} Project name (required)
            \item \textbf{Team Members:} All member names (comma-separated, required)
            \item \textbf{Supervisor:} Faculty supervisor (select from list)
            \item \textbf{Academic Year:} Year of completion (e.g., 2024, required)
            \item \textbf{Course Code:} Related course (e.g., CSE400)
            \item \textbf{Abstract:} Project summary (required)
            \item \textbf{Keywords:} Tags for searchability
        \end{itemize}
    \item Add project resources:
        \begin{itemize}
            \item Upload project file (PDF, ZIP, APK)
            \item Web URL (for live demo)
            \item GitHub repository link
            \item App download link
        \end{itemize}
    \item Choose status:
        \begin{itemize}
            \item \textbf{Draft:} Save for later
            \item \textbf{Submit for Approval:} Send to librarian/admin
        \end{itemize}
    \item Click \textbf{"Submit Project"}
    \item Your project appears in \textbf{"My Uploads"} with status "Review"
\end{enumerate}

\newpage

\subsection{Project Approval Process}
\begin{enumerate}[leftmargin=*]
    \item After submission, librarian/administrator reviews your project
    \item Review criteria:
        \begin{itemize}
            \item Completeness of information
            \item Appropriate content
            \item Working links and files
        \end{itemize}
    \item Project status changes:
        \begin{itemize}
            \item \textbf{Approved → Published:} Visible in gallery
            \item \textbf{Rejected → Draft:} Returned with feedback
        \end{itemize}
    \item You receive email notification of the decision
    \item If rejected, make requested changes and resubmit
\end{enumerate}

\subsection{Viewing Project Details}
\begin{enumerate}[leftmargin=*]
    \item Click on any project in the gallery
    \item View complete information:
        \begin{itemize}
            \item Full description and abstract
            \item Team member list
            \item Supervisor information
            \item Academic year and course
            \item Keywords and tags
        \end{itemize}
    \item Access project resources:
        \begin{itemize}
            \item Download documentation
            \item Visit web demo
            \item View source code on GitHub
            \item Download APK (for apps)
        \end{itemize}
\end{enumerate}

\section{AI-Powered Features}

\subsection{AI Chat Assistant (Student, Researcher, Librarian, Administrator)}

\subsubsection{Overview}
The AI Chat Assistant uses Retrieval-Augmented Generation (RAG) technology to answer questions by searching through all indexed documents in the platform (research papers, archives, projects).

\subsubsection{Accessing the AI Chat}
\begin{enumerate}[leftmargin=*]
    \item Click the \textbf{AI Chat icon} (usually in bottom-right corner)
    \item Or navigate to \textbf{"AI Assistant"} in the menu
    \item The chat interface opens with a welcome message
\end{enumerate}

\subsubsection{Asking Questions}
\begin{enumerate}[leftmargin=*]
    \item Type your question in the chat input box
    \item Questions can be in:
        \begin{itemize}
            \item English
            \item Bengali
            \item Other supported languages (auto-detected)
        \end{itemize}
    \item Example questions:
        \begin{itemize}
            \item "What research papers discuss machine learning algorithms?"
            \item "Summarize projects related to Android development"
            \item "Find documents about computer networking"
            \item "কম্পিউটার সায়েন্স সম্পর্কে কি আছে?" (Bengali)
        \end{itemize}
    \item Press Enter or click the Send button
    \item Wait for the AI to process (usually 2-5 seconds)
\end{enumerate}

\subsubsection{Understanding AI Responses}
AI responses include:
\begin{itemize}[leftmargin=*]
    \item \textbf{Answer:} Generated response based on indexed documents
    \item \textbf{Source Citations:} Links to documents used to generate the answer
    \item \textbf{Model Used:} Which AI model processed your query (Groq/Gemini)
    \item \textbf{Detected Language:} The language of your question
    \item \textbf{Response Time:} How long it took to generate
\end{itemize}

\begin{tcolorbox}[colback=blue!5!white,colframe=primarycolor,title=AI Capabilities]
\begin{itemize}[leftmargin=*]
    \item Searches across all published content you have access to
    \item Respects access control (won't show restricted content)
    \item Provides cited sources for transparency
    \item Remembers context within a conversation session
    \item Can rephrase questions for better results
\end{itemize}
\end{tcolorbox}

\newpage

\subsubsection{Viewing Source Documents}
\begin{enumerate}[leftmargin=*]
    \item Click on any citation link in the AI response
    \item You're taken to the source document details page
    \item Read the full document or download it
    \item Return to chat to continue conversation
\end{enumerate}

\subsubsection{Chat History}
\begin{enumerate}[leftmargin=*]
    \item Each chat creates a session
    \item View past conversations by clicking \textbf{"Chat History"}
    \item Select a previous session to continue that conversation
    \item Chat sessions are saved indefinitely
    \item Start a new session by clicking \textbf{"New Chat"}
\end{enumerate}

\subsection{Document Summarization}

\subsubsection{Generating Summaries}
\begin{enumerate}[leftmargin=*]
    \item Navigate to any research paper or document
    \item Click the \textbf{"Summarize"} button
    \item Select summary options:
        \begin{itemize}
            \item Language (English, Bengali, etc.)
            \item Model (Groq or Gemini)
            \item Length (Short, Medium, Long)
        \end{itemize}
    \item Click \textbf{"Generate Summary"}
    \item Wait for processing (10-30 seconds depending on document length)
    \item Read the AI-generated summary
    \item Copy or share the summary
\end{enumerate}

\begin{tcolorbox}[colback=yellow!10!white,colframe=orange,title=Summary Features]
\begin{itemize}[leftmargin=*]
    \item Extracts key points from long documents
    \item Maintains technical accuracy
    \item Available in multiple languages
    \item Saves time on literature review
\end{itemize}
\end{tcolorbox}

\section{Role Request System}

\subsection{Requesting a Role Upgrade (Student, Researcher)}

\subsubsection{When to Request}
Users can request role upgrades when they need additional permissions:
\begin{itemize}[leftmargin=*]
    \item \textbf{Student → Researcher:} To submit and review research papers
    \item \textbf{Student/Researcher → Librarian:} To manage library catalog and approve content
\end{itemize}

\textbf{Note:} Administrator role cannot be self-requested and must be assigned by existing administrators.

\subsubsection{Submitting a Request}
\begin{enumerate}[leftmargin=*]
    \item Click on your profile icon
    \item Select \textbf{"Request Role Upgrade"}
    \item Or navigate to \textbf{"/role-request"}
    \item Fill in the request form:
        \begin{itemize}
            \item \textbf{Desired Role:} Select Researcher or Librarian
            \item \textbf{Justification:} Explain why you need this role (required)
            \item Example justifications:
                \begin{itemize}
                    \item "I am a PhD student in the CSE department (ID: 2021-XX-XXX) and need to submit research papers"
                    \item "I am a faculty member and want to help manage the library catalog"
                    \item "I am working as a research assistant under Prof. X"
                \end{itemize}
        \end{itemize}
    \item Click \textbf{"Submit Request"}
    \item All administrators receive email notifications
    \item Your request status appears on the page
\end{enumerate}

\subsubsection{Request Status}
\begin{itemize}[leftmargin=*]
    \item \textbf{Pending:} Waiting for administrator review (you cannot submit another request while one is pending)
    \item \textbf{Approved:} Your role has been upgraded - log out and back in to see new permissions
    \item \textbf{Declined:} Request was not approved - see administrator notes for reason
\end{itemize}

\newpage

\subsubsection{Viewing Your Request History}
\begin{enumerate}[leftmargin=*]
    \item Go to the role request page
    \item Scroll down to see \textbf{"Your Requests"} section
    \item View all past requests with:
        \begin{itemize}
            \item Requested role
            \item Submission date
            \item Status (Pending, Approved, Declined)
            \item Your justification
            \item Administrator decision notes (if available)
        \end{itemize}
\end{enumerate}

\section{Notifications}

\subsection{In-App Notifications}

\subsubsection{Accessing Notifications}
\begin{enumerate}[leftmargin=*]
    \item Look for the \textbf{bell icon} in the top navigation bar
    \item Red badge shows number of unread notifications
    \item Click the bell to open notification dropdown
    \item See recent notifications with:
        \begin{itemize}
            \item Notification title
            \item Brief message
            \item Timestamp
            \item Read/unread status
        \end{itemize}
\end{enumerate}

\subsubsection{Notification Types}
You receive notifications for:
\begin{itemize}[leftmargin=*]
    \item \textbf{Role Requests:}
        \begin{itemize}
            \item Your request was approved/declined
            \item (Admins) New role request submitted
        \end{itemize}
    \item \textbf{Library:}
        \begin{itemize}
            \item Book hold is now available for pickup
            \item Loan due date approaching
            \item Overdue book reminder
        \end{itemize}
    \item \textbf{Research:}
        \begin{itemize}
            \item Paper review completed
            \item Paper published
        \end{itemize}
    \item \textbf{Projects:}
        \begin{itemize}
            \item Project approved/rejected
        \end{itemize}
    \item \textbf{System:}
        \begin{itemize}
            \item Platform updates
            \item Maintenance announcements
        \end{itemize}
\end{itemize}

\subsubsection{Managing Notifications}
\begin{enumerate}[leftmargin=*]
    \item Click on any notification to:
        \begin{itemize}
            \item Mark it as read
            \item Navigate to related page
        \end{itemize}
    \item Click \textbf{"Mark all as read"} to clear all unread notifications
    \item Click \textbf{"View All"} to see complete notification history
\end{enumerate}

\subsection{Email Notifications}
All in-app notifications are also sent to your registered email address. Configure email preferences in Settings.

\newpage

\section{Librarian Features}

\subsection{Overview (Librarian, Administrator Only)}
Librarians have additional permissions to manage the library catalog, process loans, handle fines, and approve content submissions.

\subsection{Managing Library Catalog}

\subsubsection{Adding Individual Books}
\begin{enumerate}[leftmargin=*]
    \item Navigate to \textbf{"Library"} → \textbf{"Catalog"}
    \item Click \textbf{"Add Book"} button
    \item Fill in book details:
        \begin{itemize}
            \item \textbf{Title:} Full book title (required)
            \item \textbf{Author:} Author name(s) (required)
            \item \textbf{ISBN:} ISBN-10 or ISBN-13 (optional but recommended)
            \item \textbf{Topic:} Subject area (auto-detected from title or manually selected)
            \item \textbf{Format:} Book, E-book, Reference
            \item \textbf{Year:} Publication year
            \item \textbf{Publisher:} Publishing company
            \item \textbf{Total Copies:} Number of copies available
            \item \textbf{Location:} Physical shelf location
            \item \textbf{Cover Image:} Upload book cover (optional)
        \end{itemize}
    \item Click \textbf{"Add Book"}
    \item Book is immediately added to the catalog
    \item Available copies = Total copies initially
\end{enumerate}

\subsubsection{Bulk Import from CSV}
\begin{enumerate}[leftmargin=*]
    \item Navigate to \textbf{"Admin"} → \textbf{"Catalog Management"}
    \item Click \textbf{"Import CSV"}
    \item Download the CSV template
    \item Fill in the template with book data
    \item CSV columns:
        \begin{itemize}
            \item title, author, isbn, topic, format
            \item year, publisher, total\_copies, location
        \end{itemize}
    \item Upload your completed CSV file
    \item Click \textbf{"Import"}
    \item System processes the file:
        \begin{itemize}
            \item Existing books (matched by ISBN) are updated
            \item New books are created
        \end{itemize}
    \item View import summary (successes, errors, duplicates)
\end{enumerate}

\subsubsection{Exporting Catalog}
\begin{enumerate}[leftmargin=*]
    \item Navigate to \textbf{"Admin"} → \textbf{"Catalog Management"}
    \item Click \textbf{"Export Catalog"}
    \item Select export format (CSV)
    \item Choose filters (optional):
        \begin{itemize}
            \item By topic
            \item By availability
            \item By format
        \end{itemize}
    \item Click \textbf{"Export"}
    \item CSV file downloads to your device
    \item Use for backup, analysis, or sharing
\end{enumerate}

\subsection{Managing Loans and Returns}

\subsubsection{Viewing All Loans}
\begin{enumerate}[leftmargin=*]
    \item Navigate to \textbf{"Library"} → \textbf{"All Loans"}
    \item See complete loan list with:
        \begin{itemize}
            \item Borrower name and email
            \item Book title
            \item Checkout date
            \item Due date
            \item Status (Active, Returned, Overdue)
        \end{itemize}
    \item Filter by status:
        \begin{itemize}
            \item Active loans
            \item Overdue loans
            \item Returned loans
        \end{itemize}
    \item Sort by due date, checkout date, or borrower name
\end{enumerate}

\subsubsection{Processing Returns}
\begin{enumerate}[leftmargin=*]
    \item In the All Loans list, find the loan to process
    \item Click \textbf{"Process Return"}
    \item Confirm the return
    \item System automatically:
        \begin{itemize}
            \item Marks loan as returned
            \item Calculates fines if overdue
            \item Updates book availability
            \item Notifies next user in hold queue (if any)
        \end{itemize}
\end{enumerate}

\newpage

\subsection{Fine Management}

\subsubsection{Viewing All Fines}
\begin{enumerate}[leftmargin=*]
    \item Navigate to \textbf{"Library"} → \textbf{"All Fines"}
    \item View complete fine list showing:
        \begin{itemize}
            \item Borrower name and email
            \item Book title
            \item Fine amount (BDT)
            \item Status (Pending, Paid, Waived)
            \item Calculation date
        \end{itemize}
    \item See total unpaid fines at the top
    \item Filter by status
    \item Search by borrower name
\end{enumerate}

\subsubsection{Recording Fine Payments}
\begin{enumerate}[leftmargin=*]
    \item Find the fine in the All Fines list
    \item Click \textbf{"Record Payment"}
    \item Confirm payment details:
        \begin{itemize}
            \item Amount paid
            \item Payment method
            \item Reference number (if applicable)
        \end{itemize}
    \item Click \textbf{"Confirm"}
    \item Fine status changes to "Paid"
    \item Borrower is notified
\end{enumerate}

\subsubsection{Waiving Fines}
\begin{enumerate}[leftmargin=*]
    \item Find the fine to waive
    \item Click \textbf{"Waive Fine"}
    \item Enter waiver reason (required):
        \begin{itemize}
            \item "First-time offence"
            \item "System error"
            \item "Special circumstances"
            \item Other explanation
        \end{itemize}
    \item Click \textbf{"Waive"}
    \item Fine status changes to "Waived"
    \item Action is logged in audit trail
    \item Borrower is notified
\end{enumerate}

\subsection{Content Approval}

\subsubsection{Reviewing Submitted Content}
\begin{enumerate}[leftmargin=*]
    \item Navigate to \textbf{"Admin"} → \textbf{"Content Review Queue"}
    \item See all items pending approval:
        \begin{itemize}
            \item Research papers awaiting review
            \item Student projects for approval
            \item Archive items for publishing
        \end{itemize}
    \item Click on any item to view details
    \item Review content for:
        \begin{itemize}
            \item Appropriate content
            \item Complete metadata
            \item Proper formatting
            \item Working files/links
        \end{itemize}
\end{enumerate}

\subsubsection{Approving/Rejecting Content}
\begin{enumerate}[leftmargin=*]
    \item After reviewing, choose action:
    \item \textbf{To Approve:}
        \begin{itemize}
            \item Click \textbf{"Approve"} or \textbf{"Publish"}
            \item Add approval notes (optional)
            \item Confirm
            \item Content becomes publicly visible
            \item Uploader is notified
        \end{itemize}
    \item \textbf{To Reject:}
        \begin{itemize}
            \item Click \textbf{"Reject"} or \textbf{"Request Changes"}
            \item Enter reason for rejection (required)
            \item Suggest improvements
            \item Confirm
            \item Content returns to draft status
            \item Uploader is notified with your feedback
        \end{itemize}
\end{enumerate}

\newpage

\section{Administrator Features}

\subsection{Overview (Administrator Only)}
Administrators have full platform control, including user management, role assignment, and system configuration.

\subsection{User Management}

\subsubsection{Viewing All Users}
\begin{enumerate}[leftmargin=*]
    \item Navigate to \textbf{"Admin"} → \textbf{"Users"}
    \item See complete user list with pagination
    \item For each user, view:
        \begin{itemize}
            \item Name and email
            \item Current role
            \item Registration date
            \item Last login time
        \end{itemize}
    \item Search users by name or email
    \item Filter by role
\end{enumerate}

\subsubsection{Changing User Roles}
\begin{enumerate}[leftmargin=*]
    \item Find the user in the user list
    \item Click on their row to expand details
    \item Click \textbf{"Change Role"}
    \item Select new role:
        \begin{itemize}
            \item Student
            \item Researcher
            \item Librarian
            \item Administrator
        \end{itemize}
    \item Add reason for change (recorded in audit log)
    \item Click \textbf{"Confirm"}
    \item User is notified of role change
    \item Change is logged in audit trail
\end{enumerate}

\begin{tcolorbox}[colback=red!5!white,colframe=red,title=Role Change Restrictions]
\begin{itemize}[leftmargin=*]
    \item You cannot change your own role (prevents accidental self-demotion)
    \item Administrator role should be granted carefully
    \item Role downgrades may remove user's access to content they uploaded
\end{itemize}
\end{tcolorbox}

\subsection{Role Request Management}

\subsubsection{Viewing Pending Requests}
\begin{enumerate}[leftmargin=*]
    \item Navigate to \textbf{"Admin"} → \textbf{"Role Requests"}
    \item See queue of pending requests
    \item For each request, view:
        \begin{itemize}
            \item Applicant name and email
            \item Current role
            \item Requested role
            \item Justification provided
            \item Submission date
        \end{itemize}
    \item Sort by date (oldest first recommended)
\end{enumerate}

\subsubsection{Approving Requests}
\begin{enumerate}[leftmargin=*]
    \item Click on a request to view full details
    \item Review the justification
    \item Verify applicant credentials (if needed)
    \item Click \textbf{"Approve"}
    \item Add approval notes (optional)
    \item Confirm approval
    \item System immediately:
        \begin{itemize}
            \item Updates user's role
            \item Sends email notification to user
            \item Creates audit log entry
            \item Marks request as "Approved"
        \end{itemize}
    \item User must log out and back in to see new permissions
\end{enumerate}

\subsubsection{Declining Requests}
\begin{enumerate}[leftmargin=*]
    \item Click on the request
    \item Click \textbf{"Decline"}
    \item Enter reason for decline (required):
        \begin{itemize}
            \item "Insufficient justification"
            \item "Role not appropriate for user"
            \item "Additional information needed"
            \item Other specific reason
        \end{itemize}
    \item Click \textbf{"Confirm"}
    \item User receives email with your reason
    \item Request is marked as "Declined"
    \item User can submit a new request with updated justification
\end{enumerate}

\newpage

\subsection{Audit Log}

\subsubsection{Viewing the Audit Trail}
\begin{enumerate}[leftmargin=*]
    \item Navigate to \textbf{"Admin"} → \textbf{"Audit Log"}
    \item See chronological list of all administrative actions
    \item Each entry shows:
        \begin{itemize}
            \item Actor (admin who performed action)
            \item Action type (e.g., "role\_change", "fine\_waived")
            \item Resource type (user, book, fine, etc.)
            \item Resource ID
            \item Timestamp
            \item Additional details
        \end{itemize}
    \item Filter by:
        \begin{itemize}
            \item Actor (specific admin)
            \item Action type
            \item Date range
            \item Resource type
        \end{itemize}
    \item Search by resource ID
\end{enumerate}

\subsubsection{Audit Log Actions Tracked}
\begin{itemize}[leftmargin=*]
    \item User role changes
    \item Fine waivers
    \item Role request approvals/declines
    \item Content status changes (publish, archive)
    \item Catalog imports
    \item User account modifications
    \item System configuration changes
\end{itemize}

\begin{tcolorbox}[colback=blue!5!white,colframe=primarycolor,title=Audit Log Purpose]
The audit log provides:
\begin{itemize}[leftmargin=*]
    \item Accountability for administrative actions
    \item Security and compliance monitoring
    \item Troubleshooting support
    \item Historical record of platform changes
    \item Immutable, append-only log
\end{itemize}
\end{tcolorbox}

\subsection{System Configuration}

\subsubsection{Platform Settings}
Administrators can configure:
\begin{itemize}[leftmargin=*]
    \item \textbf{Library Settings:}
        \begin{itemize}
            \item Default loan period (currently 14 days)
            \item Fine rate (currently 50 BDT/day)
            \item Maximum fine per loan (currently 500 BDT)
            \item Fine block threshold (currently 50 BDT)
        \end{itemize}
    \item \textbf{Upload Settings:}
        \begin{itemize}
            \item Maximum file size (currently 50 MB)
            \item Allowed file types
            \item Default access tiers
        \end{itemize}
    \item \textbf{AI Settings:}
        \begin{itemize}
            \item Preferred AI model
            \item Query limits
            \item Language preferences
        \end{itemize}
    \item \textbf{Email Notifications:}
        \begin{itemize}
            \item SMTP configuration
            \item Email templates
            \item Notification preferences
        \end{itemize}
\end{itemize}

\textbf{Note:} System configuration changes require environment variable updates and may require platform restart.

\newpage

\section{Search and Discovery}

\subsection{Global Search}
\begin{enumerate}[leftmargin=*]
    \item Use the search bar in the top navigation
    \item Enter keywords
    \item Search across:
        \begin{itemize}
            \item Library catalog (titles, authors, ISBN)
            \item Research papers (titles, abstracts, authors)
            \item Student projects (titles, descriptions, team members)
            \item Digital archives (titles, content, keywords)
        \end{itemize}
    \item Results are grouped by type
    \item Filter results by:
        \begin{itemize}
            \item Content type
            \item Date range
            \item Access tier
        \end{itemize}
\end{enumerate}

\subsection{Advanced Search}
\begin{enumerate}[leftmargin=*]
    \item Click \textbf{"Advanced Search"}
    \item Set multiple criteria:
        \begin{itemize}
            \item Exact phrase matching
            \item Author/creator name
            \item Date range
            \item Specific fields (title, abstract, keywords)
            \item Boolean operators (AND, OR, NOT)
        \end{itemize}
    \item Select content types to search
    \item Click \textbf{"Search"}
    \item Refine results with faceted filters
\end{enumerate}

\subsection{Browsing by Category}
\begin{enumerate}[leftmargin=*]
    \item Navigate to specific sections:
        \begin{itemize}
            \item \textbf{Library:} Browse by topic, format, or year
            \item \textbf{Research:} Browse by research area or year
            \item \textbf{Projects:} Browse by academic year or course
            \item \textbf{Archives:} Browse by collection or media type
        \end{itemize}
    \item Use hierarchical filters to narrow results
    \item View item counts for each category
\end{enumerate}

\section{Mobile Access}

\subsection{Responsive Design}
The platform is fully responsive and optimized for mobile devices:
\begin{itemize}[leftmargin=*]
    \item Adaptive layout for phones and tablets
    \item Touch-optimized interface
    \item Mobile-friendly navigation menu
    \item Optimized file uploads
    \item Full feature parity with desktop
\end{itemize}

\subsection{Mobile Browser Recommendations}
\begin{itemize}[leftmargin=*]
    \item Chrome for Android (latest version)
    \item Safari for iOS (latest version)
    \item Firefox Mobile (latest version)
\end{itemize}

\subsection{Mobile-Specific Features}
\begin{itemize}[leftmargin=*]
    \item Camera integration for document scanning
    \item Location-based library search
    \item Offline content caching (when implemented)
    \item Push notifications (when implemented)
\end{itemize}

\newpage

\section{Troubleshooting and FAQs}

\subsection{Common Issues and Solutions}

\subsubsection{Login Problems}
\begin{itemize}[leftmargin=*]
    \item \textbf{Forgot password:}
        \begin{itemize}
            \item Click "Forgot Password?" on login page
            \item Check email for reset link (check spam folder)
            \item Reset link expires in 1 hour
        \end{itemize}
    \item \textbf{Account locked after failed attempts:}
        \begin{itemize}
            \item After 5 failed login attempts, account is locked for 15 minutes
            \item Wait for the lockout period to expire
            \item Or use password reset to unlock immediately
        \end{itemize}
    \item \textbf{Email not registered:}
        \begin{itemize}
            \item Verify you're using the correct email
            \item Create a new account if needed
        \end{itemize}
\end{itemize}

\subsubsection{Upload Issues}
\begin{itemize}[leftmargin=*]
    \item \textbf{File too large:}
        \begin{itemize}
            \item Maximum file size is 50 MB
            \item Compress PDF files or videos before uploading
            \item Consider using external links for large videos
        \end{itemize}
    \item \textbf{Unsupported file type:}
        \begin{itemize}
            \item Check supported formats list
            \item Convert file to supported format
        \end{itemize}
    \item \textbf{Upload timeout:}
        \begin{itemize}
            \item Check your internet connection
            \item Try uploading during off-peak hours
            \item Upload smaller files
        \end{itemize}
\end{itemize}

\subsubsection{Borrowing Issues}
\begin{itemize}[leftmargin=*]
    \item \textbf{Cannot borrow (button disabled):}
        \begin{itemize}
            \item Check if you have outstanding fines exceeding 50 BDT
            \item Verify all copies aren't already borrowed
            \item Ensure you're logged in
            \item Librarians cannot borrow books
        \end{itemize}
    \item \textbf{Book not available:}
        \begin{itemize}
            \item Place a hold to be notified when available
            \item Check estimated return date
        \end{itemize}
\end{itemize}

\subsubsection{Search Not Working}
\begin{itemize}[leftmargin=*]
    \item \textbf{No results found:}
        \begin{itemize}
            \item Check spelling of search terms
            \item Try broader keywords
            \item Remove filters
            \item Use advanced search for more control
        \end{itemize}
    \item \textbf{Incomplete results:}
        \begin{itemize}
            \item Ensure you have permission to view restricted content
            \item Check if content is published (not draft)
        \end{itemize}
\end{itemize}

\subsubsection{AI Chat Issues}
\begin{itemize}[leftmargin=*]
    \item \textbf{AI not responding:}
        \begin{itemize}
            \item Wait up to 30 seconds for response
            \item Check your internet connection
            \item Try rephrasing your question
            \item Start a new chat session
        \end{itemize}
    \item \textbf{Irrelevant answers:}
        \begin{itemize}
            \item Be more specific in your question
            \item Include context or keywords
            \item Try asking in a different language
        \end{itemize}
\end{itemize}

\newpage

\subsection{Frequently Asked Questions}

\subsubsection{General}
\begin{enumerate}[leftmargin=*]
    \item \textbf{Q: Can I use the platform without an account?}\\
    A: Yes, you can browse published content (library catalog, research papers, projects) without logging in. However, borrowing books, uploading content, and using AI features require an account.
    
    \item \textbf{Q: How do I get elevated permissions (Researcher/Librarian)?}\\
    A: Use the Role Request feature from your profile. Provide a clear justification. An administrator will review and approve/decline your request.
    
    \item \textbf{Q: Can I delete my uploaded content?}\\
    A: You can delete draft items. Published content requires administrator approval for deletion to maintain platform integrity.
    
    \item \textbf{Q: How long does content review take?}\\
    A: Review times vary depending on administrator availability. Typically 1-5 business days. You'll receive email notification when review is complete.
\end{enumerate}

\subsubsection{Library}
\begin{enumerate}[leftmargin=*]
    \item \textbf{Q: How many books can I borrow at once?}\\
    A: There is no hard limit on simultaneous loans, but you must not have overdue fines exceeding 50 BDT.
    
    \item \textbf{Q: Can I renew a loan?}\\
    A: Currently, the system does not support automatic renewals. Return the book and re-borrow if available, or contact a librarian for assistance.
    
    \item \textbf{Q: What if I lose a borrowed book?}\\
    A: Contact the library immediately. Replacement cost will be assessed by the librarian.
    
    \item \textbf{Q: Can fines be waived?}\\
    A: Librarians and administrators can waive fines on a case-by-case basis with proper justification.
\end{enumerate}

\subsubsection{Research and Projects}
\begin{enumerate}[leftmargin=*]
    \item \textbf{Q: Who can submit research papers?}\\
    A: Only users with Researcher, Librarian, or Administrator roles can submit research papers.
    
    \item \textbf{Q: Can I edit my research paper after submission?}\\
    A: You can edit papers in draft status. After submission for review or publication, contact an administrator for changes.
    
    \item \textbf{Q: Can multiple students collaborate on a project?}\\
    A: Yes, list all team members in the project submission. The submitter becomes the primary contact.
    
    \item \textbf{Q: How do I add a supervisor to my project?}\\
    A: Select the supervisor from the user list during project submission. The supervisor must already have an account.
\end{enumerate}

\subsubsection{AI Features}
\begin{enumerate}[leftmargin=*]
    \item \textbf{Q: What content does the AI search?}\\
    A: The AI searches all published content you have permission to access, including research papers, projects, archives, and library catalog descriptions.
    
    \item \textbf{Q: Can I trust AI responses?}\\
    A: AI responses are based on indexed documents. Always verify critical information by checking the cited sources. Use AI as a starting point, not the final authority.
    
    \item \textbf{Q: What languages does the AI support?}\\
    A: The AI supports English and Bengali, with auto-detection. Additional languages may be available depending on configuration.
    
    \item \textbf{Q: Are my chat conversations private?}\\
    A: Chat sessions are associated with your account. Administrators do not routinely view chat logs, but conversations are stored for system improvement and support purposes.
\end{enumerate}

\newpage

\section{Best Practices}

\subsection{For All Users}
\begin{itemize}[leftmargin=*]
    \item \textbf{Use strong passwords:} Include uppercase, lowercase, numbers, and special characters
    \item \textbf{Keep your profile updated:} Ensure your email is current for notifications
    \item \textbf{Provide descriptive titles:} Use clear, searchable titles for uploads
    \item \textbf{Add keywords:} Enhance discoverability with relevant tags
    \item \textbf{Respect copyright:} Only upload content you have permission to share
    \item \textbf{Use appropriate access tiers:} Don't restrict public content unnecessarily
    \item \textbf{Cite sources:} When uploading research, properly attribute references
\end{itemize}

\subsection{For Students}
\begin{itemize}[leftmargin=*]
    \item \textbf{Return books on time:} Avoid fines and help others access materials
    \item \textbf{Place holds early:} Don't wait until you urgently need a book
    \item \textbf{Save drafts:} Work on projects incrementally before final submission
    \item \textbf{List all team members:} Ensure proper credit for collaborative work
    \item \textbf{Provide detailed abstracts:} Help others understand your project
    \item \textbf{Use meaningful project titles:} Avoid vague names like "Final Project"
\end{itemize}

\subsection{For Researchers}
\begin{itemize}[leftmargin=*]
    \item \textbf{Complete metadata:} Provide full author lists, keywords, and abstracts
    \item \textbf{Include DOI:} Add DOI when available for easier citation
    \item \textbf{Peer review thoughtfully:} Provide constructive feedback
    \item \textbf{Keep papers updated:} Submit revisions when significant changes occur
    \item \textbf{Collaborate openly:} Share preliminary findings when appropriate
\end{itemize}

\subsection{For Librarians}
\begin{itemize}[leftmargin=*]
    \item \textbf{Maintain catalog accuracy:} Double-check ISBN and publication details
    \item \textbf{Process returns promptly:} Keep availability information current
    \item \textbf{Document fine waivers:} Provide clear reasons in audit trail
    \item \textbf{Review content fairly:} Apply consistent approval standards
    \item \textbf{Respond to holds quickly:} Notify users when reserved books are available
\end{itemize}

\subsection{For Administrators}
\begin{itemize}[leftmargin=*]
    \item \textbf{Review role requests promptly:} Reduce waiting times for users
    \item \textbf{Document role changes:} Always include justification in audit log
    \item \textbf{Monitor system health:} Regularly check logs for errors
    \item \textbf{Backup data regularly:} Export catalogs and content lists periodically
    \item \textbf{Communicate changes:} Inform users of platform updates
\end{itemize}

\section{Security and Privacy}

\subsection{Data Protection}
\begin{itemize}[leftmargin=*]
    \item All passwords are encrypted using industry-standard bcrypt hashing
    \item Sensitive data is transmitted over HTTPS
    \item Access control enforced at database and application levels
    \item Audit logs track all administrative actions
    \item Regular security updates and patches applied
\end{itemize}

\subsection{Privacy Policy Highlights}
\begin{itemize}[leftmargin=*]
    \item Personal information (name, email) used only for platform functionality
    \item No data sharing with third parties without consent
    \item Email addresses not used for marketing
    \item User content owned by the uploader
    \item Account deletion available upon request (contact administrator)
    \item Chat logs stored for system improvement but not routinely monitored
\end{itemize}

\subsection{User Responsibilities}
\begin{itemize}[leftmargin=*]
    \item Keep your password confidential
    \item Do not share your account credentials
    \item Log out from shared computers
    \item Report suspicious activity to administrators
    \item Respect other users' privacy and content
    \item Do not attempt to access restricted content
    \item Report security vulnerabilities responsibly
\end{itemize}

\newpage

\section{Support and Contact}

\subsection{Getting Help}

\subsubsection{In-Platform Support}
\begin{itemize}[leftmargin=*]
    \item Use the \textbf{Help} link in the footer for FAQs
    \item Check this user manual for detailed instructions
    \item Use the AI Chat assistant for quick questions
\end{itemize}

\subsubsection{Contacting Support}
For technical issues or questions not covered in this manual:
\begin{itemize}[leftmargin=*]
    \item \textbf{Email:} Contact the CSE Department IT support team
    \item \textbf{Phone:} Call the department office during business hours
    \item \textbf{In-Person:} Visit the department office or library desk
\end{itemize}

\subsection{Reporting Issues}

\subsubsection{Bug Reports}
When reporting a bug, include:
\begin{itemize}[leftmargin=*]
    \item Your username (email)
    \item Browser and version
    \item Steps to reproduce the issue
    \item Screenshots if applicable
    \item Error messages
    \item Date and time of occurrence
\end{itemize}

\subsubsection{Feature Requests}
Suggest new features by contacting administrators with:
\begin{itemize}[leftmargin=*]
    \item Clear description of the proposed feature
    \item Use case and benefits
    \item Priority level (nice-to-have vs. critical)
\end{itemize}

\subsection{Platform Updates}

\subsubsection{Release Notes}
\begin{itemize}[leftmargin=*]
    \item Check the \textbf{Announcements} section for platform updates
    \item Subscribe to email notifications for major changes
    \item Review release notes for new features and bug fixes
\end{itemize}

\subsubsection{Maintenance Windows}
\begin{itemize}[leftmargin=*]
    \item Scheduled maintenance announced 48 hours in advance
    \item Typically performed during low-usage hours
    \item Expected duration communicated in advance
    \item Emergency maintenance may occur without notice
\end{itemize}

\section{Glossary}

\begin{description}[leftmargin=2cm, style=nextline]
    \item[Access Tier] Permission level required to view content (Public, Student, Researcher, Librarian, Restricted)
    
    \item[Audit Log] Immutable record of all administrative actions on the platform
    
    \item[Citation] Reference to a source document used by the AI to generate a response
    
    \item[DOI] Digital Object Identifier - unique identifier for academic papers
    
    \item[Fine] Monetary penalty for overdue book returns (50 BDT/day, max 500 BDT)
    
    \item[Hold] Reservation on a borrowed book to be notified when it becomes available
    
    \item[ISBN] International Standard Book Number - unique identifier for books
    
    \item[Loan Period] Duration for which a book can be borrowed (14 days)
    
    \item[Metadata] Descriptive information about content (title, abstract, keywords, tags)
    
    \item[RAG] Retrieval-Augmented Generation - AI technology that searches documents to answer questions
    
    \item[Role Tier] User permission level (Public, Student, Researcher, Librarian, Administrator)
    
    \item[Session] Continuous conversation thread in AI chat
    
    \item[Status] Content publication state (Draft, Review, Published, Archived)
\end{description}

\newpage

\section{Appendices}

\subsection{Appendix A: Quick Reference Card}

\begin{tcolorbox}[colback=blue!5!white,colframe=primarycolor,title=Quick Actions by Role]

\textbf{Student}
\begin{itemize}[leftmargin=*]
    \item Browse catalog → Borrow book → Return → Pay fines
    \item Upload archive/project → Edit metadata
    \item Request role upgrade
    \item Use AI chat assistant
\end{itemize}

\textbf{Researcher}
\begin{itemize}[leftmargin=*]
    \item Everything students can do PLUS:
    \item Submit research papers
    \item Review other researchers' papers
    \item Access researcher-tier content
\end{itemize}

\textbf{Librarian}
\begin{itemize}[leftmargin=*]
    \item Everything researchers can do (except borrow) PLUS:
    \item Add books to catalog
    \item Import/export catalog
    \item View/manage all loans and fines
    \item Waive fines
    \item Approve content (projects, research, archives)
\end{itemize}

\textbf{Administrator}
\begin{itemize}[leftmargin=*]
    \item Everything librarians can do PLUS:
    \item Manage users and change roles
    \item Approve/decline role requests
    \item View audit logs
    \item Configure system settings
\end{itemize}

\end{tcolorbox}

\subsection{Appendix B: Keyboard Shortcuts}

\begin{longtable}{|l|l|}
\hline
\textbf{Action} & \textbf{Shortcut} \\
\hline
\endfirsthead
\hline
\textbf{Action} & \textbf{Shortcut} \\
\hline
\endhead
\hline
\endfoot
Open search & Ctrl + K (Cmd + K on Mac) \\
\hline
Focus search bar & / \\
\hline
Open AI chat & Ctrl + . (Cmd + . on Mac) \\
\hline
Open notifications & Ctrl + N (Cmd + N on Mac) \\
\hline
Navigate to Library & Alt + L \\
\hline
Navigate to Research & Alt + R \\
\hline
Navigate to Projects & Alt + P \\
\hline
Navigate to Archives & Alt + A \\
\hline
Open profile & Alt + U \\
\hline
Logout & Alt + X \\
\hline
\end{longtable}

\subsection{Appendix C: File Format Support}

\begin{longtable}{|l|l|l|}
\hline
\textbf{Category} & \textbf{Formats} & \textbf{Max Size} \\
\hline
\endfirsthead
\hline
\textbf{Category} & \textbf{Formats} & \textbf{Max Size} \\
\hline
\endhead
\hline
\endfoot
Documents & PDF, DOCX, DOC, PPTX, PPT, XLSX, XLS & 50 MB \\
\hline
Images & JPG, JPEG, PNG, GIF & 50 MB \\
\hline
Videos & MP4 & 50 MB \\
\hline
Audio & MP3 & 50 MB \\
\hline
Archives & ZIP & 50 MB \\
\hline
Mobile Apps & APK & 50 MB \\
\hline
External Links & URL (YouTube, websites) & N/A \\
\hline
\end{longtable}

\newpage

\subsection{Appendix D: API Endpoints (For Developers)}

\begin{tcolorbox}[colback=yellow!10!white,colframe=orange,title=For Technical Users]
This section is for developers integrating with the platform or building custom tools.
\end{tcolorbox}

\subsubsection{Authentication Endpoints}
\begin{verbatim}
POST /api/v1/auth/register       - Create account
POST /api/v1/auth/login          - Login
POST /api/v1/auth/refresh        - Refresh access token
GET  /api/v1/auth/me             - Get current user info
POST /api/v1/auth/logout         - Logout
POST /api/v1/auth/forgot-password - Request password reset
POST /api/v1/auth/reset-password  - Reset password
\end{verbatim}

\subsubsection{Library Endpoints}
\begin{verbatim}
GET  /api/v1/library/catalog          - List catalog
GET  /api/v1/library/catalog/{itemId} - Get book details
POST /api/v1/library/loans            - Borrow book
GET  /api/v1/library/loans            - List user's loans
POST /api/v1/library/loans/{id}/return - Return book
GET  /api/v1/library/fines            - List user's fines
POST /api/v1/library/fines/{id}/pay   - Pay fine
POST /api/v1/library/holds            - Place hold
GET  /api/v1/library/holds            - List holds
DELETE /api/v1/library/holds/{id}     - Cancel hold
\end{verbatim}

\subsubsection{Media Endpoints}
\begin{verbatim}
POST /api/v1/media/upload              - Upload file
GET  /api/v1/media                     - List media
GET  /api/v1/media/{itemId}            - Get media details
GET  /api/v1/media/{itemId}/download   - Download file
GET  /api/v1/media/my-uploads          - List user's uploads
PATCH /api/v1/media/{itemId}/metadata  - Edit metadata
DELETE /api/v1/media/{itemId}          - Delete media
\end{verbatim}

\subsubsection{Research Endpoints}
\begin{verbatim}
POST /api/v1/research                  - Submit paper
GET  /api/v1/research                  - List papers
GET  /api/v1/research/{paperId}        - Get paper details
PUT  /api/v1/research/{paperId}        - Update paper
POST /api/v1/research/{id}/submit      - Submit for review
POST /api/v1/research/{id}/review      - Add review
POST /api/v1/research/{id}/publish     - Publish paper
\end{verbatim}

\subsubsection{AI Endpoints}
\begin{verbatim}
POST /api/v1/ai/chat                   - Send chat query
GET  /api/v1/ai/chat/history/{sessionId} - Get chat history
POST /api/v1/ai/summarize              - Summarize document
\end{verbatim}

\subsubsection{Admin Endpoints}
\begin{verbatim}
GET  /api/v1/admin/users               - List users
PATCH /api/v1/admin/users/{id}/role    - Change user role
GET  /api/v1/admin/audit-log           - View audit log
GET  /api/v1/admin/catalog/export      - Export catalog CSV
POST /api/v1/admin/catalog/import      - Import catalog CSV
GET  /api/v1/admin/role-requests       - List role requests
POST /api/v1/admin/role-requests/{id}/decide - Approve/decline
\end{verbatim}

\newpage

\subsection{Appendix E: System Architecture Overview}

\subsubsection{Technology Stack}
\begin{itemize}[leftmargin=*]
    \item \textbf{Frontend:} Next.js 14 (React framework), TypeScript, Tailwind CSS
    \item \textbf{Backend API:} Go (Golang), Chi router
    \item \textbf{Database:} PostgreSQL 16 with pgvector extension
    \item \textbf{File Storage:} MinIO (S3-compatible object storage)
    \item \textbf{Cache/Queue:} Redis
    \item \textbf{AI/RAG Service:} Python FastAPI, Sentence Transformers
    \item \textbf{LLM Providers:} Groq (primary), Google Gemini (fallback)
    \item \textbf{Web Server:} Nginx (reverse proxy)
\end{itemize}

\subsubsection{Key Features}
\begin{itemize}[leftmargin=*]
    \item JWT-based authentication with refresh token rotation
    \item Role-based access control (RBAC) at database and application levels
    \item Vector embeddings for semantic search (768-dimensional)
    \item Hybrid search combining full-text and semantic similarity
    \item Real-time notifications via WebSocket (planned)
    \item Audit logging for compliance and security
    \item Automated fine calculation via background worker
    \item Email notifications via SMTP (Brevo/Sendinblue)
\end{itemize}

\subsection{Appendix F: Default User Accounts}

\textbf{For testing and initial setup, the following accounts are pre-configured:}

\begin{tcolorbox}[colback=red!5!white,colframe=red,title=Security Warning]
Change these default passwords immediately in production!
\end{tcolorbox}

\begin{longtable}{|l|l|l|}
\hline
\textbf{Email} & \textbf{Password} & \textbf{Role} \\
\hline
\endfirsthead
\hline
\textbf{Email} & \textbf{Password} & \textbf{Role} \\
\hline
\endhead
\hline
\endfoot
admin@cs.du.ac.bd & Admin@12345 & Administrator \\
\hline
librarian@cs.du.ac.bd & Staff@12345 & Librarian \\
\hline
researcher@cs.du.ac.bd & Research@12345 & Researcher \\
\hline
student@cs.du.ac.bd & Student@12345 & Student \\
\hline
\end{longtable}

\subsection{Appendix G: Configuration Settings}

\subsubsection{Library Configuration}
\begin{itemize}[leftmargin=*]
    \item \textbf{LOAN\_PERIOD\_DAYS:} 14 days
    \item \textbf{FINE\_RATE\_BDT\_PER\_DAY:} 50 BDT
    \item \textbf{MAX\_FINE\_PER\_LOAN\_BDT:} 500 BDT
    \item \textbf{FINE\_BLOCK\_THRESHOLD\_BDT:} 50 BDT
    \item \textbf{FINE\_CALC\_INTERVAL:} 24 hours
\end{itemize}

\subsubsection{Upload Configuration}
\begin{itemize}[leftmargin=*]
    \item \textbf{Maximum File Size:} 50 MB per file
    \item \textbf{Allowed Document Formats:} PDF, DOCX, DOC, PPTX, PPT, XLSX, XLS
    \item \textbf{Allowed Media Formats:} MP4, MP3, JPG, JPEG, PNG, GIF
    \item \textbf{Allowed Archive Formats:} ZIP, APK
\end{itemize}

\subsubsection{AI/RAG Configuration}
\begin{itemize}[leftmargin=*]
    \item \textbf{Embedding Model:} sentence-transformers/paraphrase-multilingual-MiniLM-L12-v2
    \item \textbf{Embedding Dimension:} 768
    \item \textbf{Chunk Size:} 512 tokens
    \item \textbf{Chunk Overlap:} 50 tokens
    \item \textbf{Top-K Results:} 10
    \item \textbf{Vector Search Limit:} 8
    \item \textbf{FTS Search Limit:} 8
    \item \textbf{Max Query Length:} 5000 characters
    \item \textbf{Max Context Window:} 10 previous messages
\end{itemize}

\subsubsection{Authentication Configuration}
\begin{itemize}[leftmargin=*]
    \item \textbf{JWT\_EXPIRY\_HOURS:} 1 hour
    \item \textbf{REFRESH\_EXPIRY\_DAYS:} 7 days
    \item \textbf{Bcrypt Cost:} 12 rounds
    \item \textbf{Failed Login Attempts:} 5 attempts
    \item \textbf{Account Lockout Duration:} 15 minutes
    \item \textbf{Password Reset Link Expiry:} 1 hour
\end{itemize}

\newpage

\subsection{Appendix H: Change Log}

\subsubsection{Version 1.0.0 (Initial Release)}
\textbf{Release Date:} July 2026

\textbf{Core Features:}
\begin{itemize}[leftmargin=*]
    \item User authentication and authorization system
    \item Five-tier role hierarchy (Public, Student, Researcher, Librarian, Administrator)
    \item Library catalog management with search and filters
    \item Book borrowing and return workflow
    \item Automated fine calculation (50 BDT/day, max 500 BDT)
    \item Hold/reservation system
    \item Digital archive management
    \item Research paper submission and review workflow
    \item Student project showcase
    \item Media upload (PDF, DOCX, images, videos, etc.)
    \item AI-powered chat assistant with RAG
    \item Document summarization
    \item Role request system
    \item In-app and email notifications
    \item Audit logging for administrative actions
    \item CSV import/export for library catalog
    \item Mobile-responsive design
\end{itemize}

\textbf{Security Features:}
\begin{itemize}[leftmargin=*]
    \item JWT-based authentication
    \item Refresh token rotation
    \item Bcrypt password hashing (cost 12)
    \item Brute-force protection (5 attempts, 15-min lockout)
    \item Role-based access control
    \item Audit trail for admin actions
\end{itemize}

\textbf{Known Limitations:}
\begin{itemize}[leftmargin=*]
    \item No automatic loan renewals
    \item Maximum file size limited to 50 MB
    \item No offline mode
    \item AI chat limited to English and Bengali
    \item No real-time WebSocket notifications (email only)
\end{itemize}

\subsection{Appendix I: Roadmap (Planned Features)}

\subsubsection{Upcoming Features}
\begin{itemize}[leftmargin=*]
    \item \textbf{Enhanced Search:}
        \begin{itemize}
            \item Advanced Boolean operators
            \item Saved search queries
            \item Search history
        \end{itemize}
    \item \textbf{Social Features:}
        \begin{itemize}
            \item User profiles with bios
            \item Follow other researchers
            \item Comment on research papers
            \item Rate and review books
        \end{itemize}
    \item \textbf{Collaboration:}
        \begin{itemize}
            \item Shared project workspaces
            \item Real-time document editing
            \item Discussion forums
        \end{itemize}
    \item \textbf{Analytics:}
        \begin{itemize}
            \item Usage statistics dashboard
            \item Popular content reports
            \item User engagement metrics
        \end{itemize}
    \item \textbf{Mobile App:}
        \begin{itemize}
            \item Native iOS and Android apps
            \item Offline reading
            \item Push notifications
            \item Barcode scanning for checkouts
        \end{itemize}
    \item \textbf{Integration:}
        \begin{itemize}
            \item OAuth/SSO with university systems
            \item Integration with institutional repositories
            \item API for third-party applications
        \end{itemize}
\end{itemize}

\newpage

\section{Conclusion}

The CSEDU Digital Knowledge Platform represents a comprehensive solution for managing library resources, research papers, student projects, and digital archives within an academic environment. By integrating traditional library management with modern AI-powered features, the platform facilitates knowledge sharing, collaboration, and discovery.

\subsection{Key Takeaways}

\begin{itemize}[leftmargin=*]
    \item \textbf{Unified Access:} Single platform for all departmental knowledge resources
    \item \textbf{Role-Based Permissions:} Appropriate access levels for different user types
    \item \textbf{Intelligent Search:} AI-powered discovery across all content types
    \item \textbf{Streamlined Workflows:} Automated processes for borrowing, fines, and approvals
    \item \textbf{Transparency:} Audit logs and clear status tracking
    \item \textbf{Accessibility:} Mobile-responsive design, multilingual support
\end{itemize}

\subsection{Getting the Most from the Platform}

To maximize your experience:
\begin{enumerate}[leftmargin=*]
    \item Explore all sections to understand available resources
    \item Use descriptive titles and keywords when uploading content
    \item Leverage the AI assistant for quick information retrieval
    \item Keep your profile updated for accurate notifications
    \item Provide feedback to help improve the platform
    \item Respect due dates to avoid fines and help others
    \item Collaborate with colleagues through shared content
\end{enumerate}

\subsection{Continuous Improvement}

The platform is actively maintained and regularly updated with:
\begin{itemize}[leftmargin=*]
    \item New features based on user feedback
    \item Security patches and bug fixes
    \item Performance optimizations
    \item Enhanced AI capabilities
    \item Expanded content collections
\end{itemize}

\subsection{Your Feedback Matters}

We encourage all users to:
\begin{itemize}[leftmargin=*]
    \item Report bugs and technical issues
    \item Suggest feature enhancements
    \item Share best practices with the community
    \item Contribute to documentation improvements
    \item Participate in user surveys
\end{itemize}

\vspace{1cm}

\begin{center}
\Large
\textbf{Thank you for using the CSEDU Digital Knowledge Platform!}
\end{center}

\vspace{0.5cm}

\begin{center}
For support, contact the CSE Department IT team\\
Department of Computer Science and Engineering\\
University of Dhaka\\
Dhaka-1000, Bangladesh
\end{center}

\vspace{1cm}

\begin{center}
\rule{0.5\textwidth}{0.4pt}\\
\vspace{0.3cm}
\textit{Version 1.0.0 | July 2026}\\
\textit{Document Generated: \today}
\end{center}

\end{document}
